\scanpage{73}
et $S,J$ sont deux parties discrètes disjointes de $\hat X$.

On va \deleted{\unsure{en}} passer aux catégories isotopiques correspondantes, et en donner des interprétations en termes de théorie des groupes. Bien entendu, pour cette description, on est ramené au cas connexe. On peut \deleted{aussi} noter tt de suite que pour ce pb, on est ramené au cas où $J=\varnothing$ [en \deleted{remplaçant}\inserted{\unsure{donnant|décrivant}} dans le cas général comme celui de la donnée d'un $(\hat X,I)$ \uline{plus} la donnée d'une partie $J$ de $I$]. Mais on gardera le pb sous \inserted{la} forme générale, en gardant également la dissymétrie dans le rôle des $S,J$ -- mais
% archive sheet number: 72
